\begin{filecontents*}{references.bib}
@misc{3dshapes18,
  title={3D Shapes Dataset},
  author={Chris Burgess and Hyunjik Kim},
  howpublished={\url{https://github.com/google-deepmind/3d-shapes}},
  year={2018}
}

@article{Oquab2024DINOv2,
  title={DINOv2: Learning Robust Visual Features without Supervision},
  author={Maxime Oquab and Timothee Darcet and Theo Moutakanni and Huy Vo and Marc Szafraniec and Vasil Khalidov and Pierre Fernandez and Daniel Haziza and Francisco Massa and Alaaeldin El-Nouby and Mahmoud Assran and Nicolas Ballas and Wojciech Galuba and Michael Howes and Po-Yao Huang and Ishan Misra and Quentin Deforme and Vivek Sharma and Gabriel Synnaeve and Hu Xu and Ross Girshick and Herve Jegou},
  journal={Transactions on Machine Learning Research},
  year={2024},
  url={https://arxiv.org/abs/2304.07193}
}

@article{Stevens2025Interpretable,
  title={Interpretable and Testable Vision Features via Sparse Autoencoders},
  author={Samuel Stevens and Wei-Lun Chao and Tanya Berger-Wolf and Yu Su},
  journal={arXiv preprint arXiv:2502.06755},
  year={2025},
  url={https://arxiv.org/abs/2502.06755}
}

@article{Fel2025Archetypal,
  title={Archetypal SAE: Adaptive and Stable Dictionary Learning for Concept Extraction in Large Vision Models},
  author={Thomas Fel and Ekdeep Singh Lubana and Jacob S. Prince and Matthew Kowal and Victor Boutin and Isabel Papadimitriou and Binxu Wang and Martin Wattenberg and Demba Ba and Talia Konkle},
  journal={arXiv preprint arXiv:2502.12892},
  year={2025},
  url={https://arxiv.org/abs/2502.12892}
}

@article{Costa2025From,
  title={From Flat to Hierarchical: Extracting Sparse Representations with Matching Pursuit},
  author={Valerie Costa and Thomas Fel and Ekdeep Singh Lubana and Bahareh Tolooshams and Demba Ba},
  journal={arXiv preprint arXiv:2506.03093},
  year={2025},
  url={https://arxiv.org/abs/2506.03093}
}

@article{Han2025Causal,
  title={Causal Interpretation of Sparse Autoencoder Features in Vision},
  author={Sangyu Han and Yearim Kim and Nojun Kwak},
  journal={arXiv preprint arXiv:2509.00749},
  year={2025},
  url={https://arxiv.org/abs/2509.00749}
}

@article{Ding2025ConceptSAE,
  title={Concept-SAE: Active Causal Probing of Visual Model Behavior},
  author={Jianrong Ding and Muxi Chen and Chenchen Zhao and Qiang Xu},
  journal={arXiv preprint arXiv:2509.22015},
  year={2025},
  url={https://arxiv.org/abs/2509.22015}
}

@article{Leask2025Sparse,
  title={Sparse Autoencoders Do Not Find Canonical Units of Analysis},
  author={Patrick Leask and Bart Bussmann and Michael Pearce and Joseph Bloom and Curt Tigges and Noura Al Moubayed and Lee Sharkey and Neel Nanda},
  journal={arXiv preprint arXiv:2502.04878},
  year={2025},
  url={https://arxiv.org/abs/2502.04878}
}

@article{Hindupur2025Projecting,
  title={Projecting Assumptions: The Duality Between Sparse Autoencoders and Concept Geometry},
  author={Sai Sumedh R. Hindupur and Ekdeep Singh Lubana and Thomas Fel and Demba Ba},
  journal={arXiv preprint arXiv:2503.01822},
  year={2025},
  url={https://arxiv.org/abs/2503.01822}
}

@InProceedings{Olson_2025_ICCV,
  author    = {Olson, Matthew Lyle and Hinck, Musashi and Ratzlaff, Neale and Li, Changbai and Howard, Phillip and Lal, Vasudev and Tseng, Shao-Yen},
  title     = {Probing the Representational Power of Sparse Autoencoders in Vision Models},
  booktitle = {Proceedings of the IEEE/CVF International Conference on Computer Vision (ICCV) Workshops},
  month     = {October},
  year      = {2025},
  pages     = {6226--6236},
  url       = {https://openaccess.thecvf.com/content/ICCV2025W/Findings/html/Olson_Probing_the_Representational_Power_of_Sparse_Autoencoders_in_Vision_Models_ICCVW_2025_paper.html}
}

@article{Bau2017Network,
  title={Network Dissection: Quantifying Interpretability of Deep Visual Representations},
  author={David Bau and Bolei Zhou and Aditya Khosla and Aude Oliva and Antonio Torralba},
  journal={arXiv preprint arXiv:1704.05796},
  year={2017},
  url={https://arxiv.org/abs/1704.05796}
}

@article{Kim2017Interpretability,
  title={Interpretability Beyond Feature Attribution: Quantitative Testing with Concept Activation Vectors (TCAV)},
  author={Been Kim and Martin Wattenberg and Justin Gilmer and Carrie Cai and James Wexler and Fernanda Viegas and Rory Sayres},
  journal={arXiv preprint arXiv:1711.11279},
  year={2017},
  url={https://arxiv.org/abs/1711.11279}
}

@article{Koh2020Concept,
  title={Concept Bottleneck Models},
  author={Pang Wei Koh and Thao Nguyen and Yew Siang Tang and Stephen Mussmann and Emma Pierson and Been Kim and Percy Liang},
  journal={arXiv preprint arXiv:2007.04612},
  year={2020},
  url={https://arxiv.org/abs/2007.04612}
}

@article{NasiriSarvi2025SPARC,
  title={SPARC: Concept-Aligned Sparse Autoencoders for Cross-Model and Cross-Modal Interpretability},
  author={Ali Nasiri-Sarvi and Hassan Rivaz and Mahdi S. Hosseini},
  journal={arXiv preprint arXiv:2507.06265},
  year={2025},
  url={https://arxiv.org/abs/2507.06265}
}
\end{filecontents*}

\documentclass{article}

\usepackage[preprint]{neurips_2023}
\usepackage[utf8]{inputenc}
\usepackage[T1]{fontenc}
\usepackage{hyperref}
\usepackage{url}
\usepackage{booktabs}
\usepackage{amsfonts}
\usepackage{nicefrac}
\usepackage{microtype}
\usepackage{xcolor}
\usepackage{amsmath,amssymb}
\usepackage{graphicx}
\usepackage{subcaption}

\graphicspath{{figures/}}

\newcommand{\R}{\mathbb{R}}
\newcommand{\E}{\mathbb{E}}
\newcommand{\norm}[1]{\left\lVert #1 \right\rVert}

\title{Pair-Supervised Regularization for Selective Counterfactual Edits in Frozen Vision SAEs}

\author{
  Anonymous Authors
}

\begin{document}

\maketitle

\begin{abstract}
We study a narrow question in vision interpretability: can pair-supervised regularization make sparse autoencoder (SAE) edits on frozen vision features more selective under a shared evaluation pipeline? We train width-1024 SAEs on frozen DINOv2 ViT-S/14 mean-pooled patch features and evaluate them on Shapes3D, where exact single-factor counterfactual pairs are available. Our regularizer combines counterfactual-difference concentration, same-factor signed-direction alignment, and cross-factor separation, while nuisance invariance is audit-gated rather than assumed. The audit rejects all candidate nuisance augmentation families, with worst-factor agreement ranging from 0.560 to 0.756 rather than the required 0.99, so the final claim excludes invariance. Across three seeds and two matched sparsity budgets, the pair-regularized SAE modestly increases latent concentration at budget A (0.0871 versus 0.0864 for vanilla SAE) but does not improve the primary intervention metrics over matched baselines. At budget B under rule A, vanilla SAE attains the best selective intervention score among the three main methods ($-0.1935 \pm 0.0022$) while the pair-regularized model scores $-0.1971 \pm 0.0007$. A one-seed MP-SAE baseline further improves reconstruction, off-target preservation, and retrieval consistency. The result is therefore negative but informative: in this controlled frozen-feature setting, modest pair regularization does not outperform stronger or simpler SAE baselines on edit selectivity.
\end{abstract}

\section{Introduction}

Sparse autoencoders are now a standard tool for analyzing internal representations, and recent vision work has made the space substantially more crowded. Frozen-vision SAEs already support causal editing and targeted interventions \citep{Stevens2025Interpretable,Han2025Causal}; alternative objectives can improve geometry or stability \citep{Fel2025Archetypal,Hindupur2025Projecting}; and supervised variants can impose concept alignment or grounding \citep{Ding2025ConceptSAE,NasiriSarvi2025SPARC}. In that context, a credible contribution is no longer ``SAEs make vision models interpretable.'' A narrower and testable question is whether paired supervision can improve \emph{edit selectivity}: when only one generative factor changes, does the learned latent representation support stronger target changes with fewer off-target effects?

We investigate that question in a controlled setting. The backbone is frozen DINOv2 ViT-S/14 \citep{Oquab2024DINOv2}; the benchmark is Shapes3D \citep{3dshapes18}, which supplies exact single-factor counterfactuals; and the comparison is restricted to matched-width, matched-budget SAEs trained with the same checkpoint selection rule. The proposed method augments a standard SAE objective with pair-conditioned losses that encourage sparse, repeatable latent differences for exact counterfactual pairs. To avoid overstating nuisance invariance, we first audit all candidate nuisance transformations and only keep them if factor agreement exceeds 99\% for every factor.

The main empirical result is negative. No candidate nuisance family passes the audit, so the invariance term is disabled. The remaining pair regularization slightly increases concentration diagnostics, but it does not beat both vanilla SAE and Archetypal SAE on the primary intervention metrics under either edit rule. At the higher matched sparsity budget, vanilla SAE obtains the strongest selective score among the three main methods, and a one-seed MP-SAE baseline performs better still on reconstruction, off-target preservation, and retrieval consistency.

Our contributions are:
\begin{itemize}
  \item We formulate a controlled evaluation protocol for pair-regularized vision SAEs on exact counterfactual pairs, including audited nuisance gating, matched sparsity budgets, and two shared edit rules.
  \item We show that in this setting, pair regularization modestly improves latent concentration but does not improve the main intervention metrics over vanilla SAE or Archetypal SAE.
  \item We provide a negative result that narrows the claim space for supervised vision SAEs: stronger sparse-coding choices matter more than this form of pair regularization in the evaluated frozen-feature regime.
\end{itemize}

Section~\ref{sec:related} reviews related work, Section~\ref{sec:method} describes the model and audit-gated training objective, Section~\ref{sec:experiments} presents the Shapes3D setup and results, and Section~\ref{sec:discussion} discusses limitations and implications.

\section{Related Work}
\label{sec:related}

\paragraph{Vision SAEs and causal editing.}
Standard vision SAEs already support meaningful intervention-based analysis. \citet{Stevens2025Interpretable} demonstrate editable sparse features on frozen visual backbones, and \citet{Han2025Causal} study causal interpretation of those features through image-space interventions. Our question is narrower: not whether editing is possible, but whether paired supervision during SAE training improves the \emph{selectivity} of those edits under a fixed downstream editing protocol.

\paragraph{Alternative SAE objectives.}
Archetypal SAE regularizes dictionary geometry and stability \citep{Fel2025Archetypal}, while MP-SAE changes the sparse coding family itself \citep{Costa2025From}. Theoretical and empirical work has also emphasized that SAE behavior depends strongly on data geometry and should not be equated with canonical semantic units \citep{Hindupur2025Projecting,Leask2025Sparse}. Our method differs from these objectives by supervising how latents move across exact counterfactual pairs, not by claiming a new sparse-coding family or a canonical neuron-level decomposition.

\paragraph{Supervised concept-grounded SAEs.}
Several recent papers use supervision to align SAE features to human-interpretable concepts or to enable targeted interventions \citep{Ding2025ConceptSAE,NasiriSarvi2025SPARC}. Those papers make it clear that ``adding supervision'' is not itself novel. The present study isolates a more limited question: whether pair labels alone can improve counterfactual edit behavior without concept names, masks, or segmentation labels.

\paragraph{Evaluation philosophy.}
We follow intervention-centered evaluation rather than purely qualitative feature inspection, in line with earlier concept and interpretability literature \citep{Bau2017Network,Kim2017Interpretability,Koh2020Concept}. Recent evaluation-focused vision-SAE work makes the same point: useful claims should rest on downstream behavior, not only proxy notions of monosemanticity \citep{Olson_2025_ICCV}.

\section{Method}
\label{sec:method}

\subsection{Frozen-feature SAE with pair regularization}

Let $h \in \R^d$ be a frozen DINOv2 feature for an input image. An SAE produces a sparse code $z \in \R^m$ and reconstruction $\hat{h}$ through
\[
z = f_{\theta}(h), \qquad \hat{h} = g_{\phi}(z),
\]
with width $m=1024$. The training objective starts from the usual reconstruction and sparsity terms,
\[
\mathcal{L}_{\text{base}} = \mathcal{L}_{\text{rec}} + \lambda_s \mathcal{L}_{\text{sparse}},
\]
where $\mathcal{L}_{\text{rec}} = \norm{h-\hat{h}}_2^2$ on standardized frozen features and $\mathcal{L}_{\text{sparse}} = \norm{z}_1$.

For an exact single-factor counterfactual pair $(x, x_{\mathrm{cf}}^k)$ that changes factor $k$, we define the signed latent difference
\[
d = z_{\mathrm{cf}} - z, \qquad \delta = |d|.
\]
The concentration term encourages a small subset of latents to carry most of the mass:
\[
p_i = \frac{\delta_i}{\sum_j \delta_j + \varepsilon}, \qquad
C(\delta) = \frac{m \sum_i p_i^2 - 1}{m-1},
\]
\[
\mathcal{L}_{\text{conc}} = 1 - C(\delta).
\]
Each factor $k$ also maintains a running signed centroid $\mu_k = \E[d \mid k]$. The alignment term encourages same-factor deltas to point in a repeatable direction,
\[
\mathcal{L}_{\text{align}} = 1 - \cos(d, \mu_k),
\]
and the separation term discourages different factors from collapsing to the same average direction:
\[
\mathcal{L}_{\text{sep}} = \frac{1}{K(K-1)} \sum_{k \neq k'} \max(0, \cos(\mu_k, \mu_{k'}))^2.
\]
The resulting training objective is
\[
\mathcal{L}
= \mathcal{L}_{\text{rec}}
+ \lambda_s \mathcal{L}_{\text{sparse}}
+ \lambda_c \mathcal{L}_{\text{conc}}
+ \lambda_a \mathcal{L}_{\text{align}}
+ \lambda_o \mathcal{L}_{\text{sep}}
+ \lambda_i \mathcal{L}_{\text{inv}}.
\]

\subsection{Audit-gated nuisance invariance}

The invariance term is only valid if ``nuisance-preserving'' pairs are actually nuisance-preserving in frozen feature space. We therefore audit each augmentation family using linear factor heads on Shapes3D. A family is accepted only if validation agreement exceeds 99\% for \emph{all} six factors. In the executed rerun, no family passes this threshold: the minimum factor agreement is 0.560 for crop-plus-photo perturbations, 0.629 for full photometric perturbations, and 0.756 even for the strict photometric setting. Consequently, the effective invariance weight is zero in all main runs, and the study becomes a concentration/alignment regularization experiment rather than an invariance experiment.

\subsection{Shared edit rules}

All methods are evaluated with the same two edit rules.

\paragraph{Rule A: sparse association edit.}
For each factor $k$ and latent $j$, we compute
\[
s_{k,j} = \E\!\left[|z_j(x)-z_j(x_{\mathrm{cf}})| \mid k\right]
- \E\!\left[|z_j(x)-z_j(x_{\mathrm{cf}})| \mid k' \neq k\right].
\]
We keep the top 8 latents and apply the mean signed shift estimated from training pairs on those latents only.

\paragraph{Rule B: dense centroid edit.}
We use the full signed centroid $\mu_k = \E[z_{\mathrm{cf}}-z \mid k]$ with no top-$K$ truncation.

For both rules, the edit scale $\alpha$ is selected from $\{0.5, 1.0, 1.5\}$ on the validation split only.

\section{Experiments}
\label{sec:experiments}

\subsection{Setup}

\paragraph{Data.}
Shapes3D is the only claim-bearing benchmark. We use splits of 18{,}000 train, 3{,}000 validation, and 3{,}000 test images, together with 24{,}000/4{,}000/4{,}000 exact counterfactual pairs across train/validation/test. Candidate nuisance pairs are also constructed at 18{,}000/3{,}000/3{,}000, but the nuisance audit fails and those pairs are not used to activate the invariance loss.

\paragraph{Representation selection.}
The representation pilot evaluates final-layer CLS, penultimate-layer CLS, and final-layer mean-pooled patch tokens on a reduced split. Mean-pooled patches are selected because they achieve the best selective score under rule A ($-0.254$ versus $-0.276$ for penultimate CLS and $-0.290$ for final CLS) while staying within 5\% of the best pilot reconstruction error.

\paragraph{Models and optimization.}
All main models use width 1024, batch size 1024, Adam with learning rate $10^{-3}$ and weight decay $10^{-5}$, and 18 training epochs on one RTX A6000 GPU. We evaluate two matched sparsity budgets enforced through an explicit active-set cap: budget A targets 12 active latents and budget B targets 18. Hyperparameters are tuned on seed 11 and then frozen for seeds 17 and 23. The pair model uses $\lambda_c=0.2$, $\lambda_a=0.1$, $\lambda_o=0.05$, and a tuned pair multiplier of 0.5 at budget B; the effective invariance weight is 0 after the failed audit. Checkpoints are selected using the same rule for every method: best validation reconstruction among checkpoints within the target active-latent window.

\paragraph{Metrics.}
The primary metrics are target-factor success, off-target preservation, selective intervention score, and counterfactual consistency. The selective score is
\[
\text{selective} = \text{target success} - (1 - \text{off-target preservation}),
\]
so less-negative values are better. We also report normalized reconstruction error and latent concentration. All main comparisons use three seeds; MP-SAE is a one-seed confirmatory baseline. Bootstrap confidence intervals are computed over held-out test pairs.

\subsection{Main results}

Table~\ref{tab:main} reports the central comparison. The key result is that the pair-regularized model does not beat both vanilla SAE and Archetypal SAE on the primary intervention metrics at either matched sparsity budget. At budget A, the pair model slightly increases concentration (0.0871 versus 0.0864 for vanilla and 0.0862 for archetypal), but vanilla still achieves the best selective score under both edit rules. At budget B, vanilla again performs best among the three main methods, with rule-A selective score $-0.1935 \pm 0.0022$ versus $-0.1971 \pm 0.0007$ for the pair model and $-0.1950 \pm 0.0019$ for Archetypal SAE.

The bootstrap intervals overlap heavily. For budget B rule A, the 95\% bootstrap interval for vanilla selective score is $[-0.1976,-0.1892]$, while the pair model yields $[-0.2013,-0.1928]$. The same pattern holds for consistency: vanilla reaches 0.5670 with interval $[0.5577,0.5755]$, and the pair model reaches 0.5666 with interval $[0.5576,0.5754]$. The evidence therefore does not support a claim that pair regularization improves intervention performance in this setting.

\begin{table*}[t]
\caption{Main Shapes3D comparison at matched sparsity budgets. Means and standard deviations are over seeds 11, 17, and 23, except MP-SAE which uses one seed. Best value among the three main methods in each budget/rule block is in \textbf{bold}. Lower reconstruction error is better; higher target success, off-target preservation, and consistency are better; less-negative selective score is better.}
\label{tab:main}
\centering
\small
\resizebox{\textwidth}{!}{%
\begin{tabular}{llccccc}
\toprule
Method & Budget & Norm. recon. $\downarrow$ & Target A $\uparrow$ & Off-target A $\uparrow$ & Selective A $\uparrow$ & Consistency A $\uparrow$ \\
\midrule
Vanilla SAE & A & \textbf{0.3539 $\pm$ 0.0010} & \textbf{0.0323 $\pm$ 0.0023} & \textbf{0.7261 $\pm$ 0.0041} & \textbf{-0.2417 $\pm$ 0.0032} & \textbf{0.5512 $\pm$ 0.0016} \\
Archetypal SAE & A & 0.3542 $\pm$ 0.0009 & 0.0302 $\pm$ 0.0015 & 0.7236 $\pm$ 0.0036 & -0.2463 $\pm$ 0.0022 & 0.5504 $\pm$ 0.0011 \\
Pair SAE & A & 0.3540 $\pm$ 0.0003 & 0.0313 $\pm$ 0.0020 & 0.7237 $\pm$ 0.0027 & -0.2450 $\pm$ 0.0015 & 0.5509 $\pm$ 0.0014 \\
\midrule
Vanilla SAE & B & \textbf{0.3262 $\pm$ 0.0002} & \textbf{0.0273 $\pm$ 0.0015} & \textbf{0.7792 $\pm$ 0.0030} & \textbf{-0.1935 $\pm$ 0.0022} & \textbf{0.5670 $\pm$ 0.0022} \\
Archetypal SAE & B & 0.3263 $\pm$ 0.0004 & 0.0266 $\pm$ 0.0018 & 0.7784 $\pm$ 0.0036 & -0.1950 $\pm$ 0.0019 & 0.5645 $\pm$ 0.0027 \\
Pair SAE & B & 0.3265 $\pm$ 0.0002 & 0.0260 $\pm$ 0.0024 & 0.7769 $\pm$ 0.0020 & -0.1971 $\pm$ 0.0007 & 0.5666 $\pm$ 0.0039 \\
\midrule
MP-SAE & B & 0.2893 & 0.0180 & 0.8415 & -0.1405 & 0.5935 \\
\bottomrule
\end{tabular}
}
\vspace{0.5em}

\resizebox{0.92\textwidth}{!}{%
\begin{tabular}{llcccc}
\toprule
Method & Budget & Target B $\uparrow$ & Off-target B $\uparrow$ & Selective B $\uparrow$ & Consistency B $\uparrow$ \\
\midrule
Vanilla SAE & A & \textbf{0.0326 $\pm$ 0.0028} & \textbf{0.7260 $\pm$ 0.0042} & \textbf{-0.2415 $\pm$ 0.0034} & 0.5506 $\pm$ 0.0014 \\
Archetypal SAE & A & 0.0304 $\pm$ 0.0013 & 0.7235 $\pm$ 0.0037 & -0.2461 $\pm$ 0.0025 & 0.5503 $\pm$ 0.0010 \\
Pair SAE & A & 0.0318 $\pm$ 0.0030 & 0.7235 $\pm$ 0.0029 & -0.2447 $\pm$ 0.0014 & \textbf{0.5512 $\pm$ 0.0014} \\
\midrule
Vanilla SAE & B & \textbf{0.0280 $\pm$ 0.0014} & \textbf{0.7794 $\pm$ 0.0030} & \textbf{-0.1927 $\pm$ 0.0027} & \textbf{0.5668 $\pm$ 0.0020} \\
Archetypal SAE & B & 0.0267 $\pm$ 0.0025 & 0.7787 $\pm$ 0.0038 & -0.1947 $\pm$ 0.0018 & 0.5648 $\pm$ 0.0026 \\
Pair SAE & B & 0.0262 $\pm$ 0.0025 & 0.7762 $\pm$ 0.0026 & -0.1975 $\pm$ 0.0003 & 0.5665 $\pm$ 0.0032 \\
\midrule
MP-SAE & B & 0.0188 & 0.8410 & -0.1402 & 0.5928 \\
\bottomrule
\end{tabular}
}
\end{table*}

Figures~\ref{fig:ruleA} and \ref{fig:ruleB} visualize the seed-level comparisons for both edit rules. Figure~\ref{fig:tradeoff} shows the reconstruction-selectivity trade-off: moving from budget A to B helps all three main methods, but the pair-regularized model never occupies the best frontier point. In this experiment, the stronger family baseline MP-SAE dominates the trade-off despite its lower target success, because it preserves non-target factors much better.

\begin{figure*}[t]
\centering
\includegraphics[width=0.9\linewidth]{main_metrics_ruleA.pdf}
\caption{Seed-level Shapes3D results under rule A, the sparse association edit. Pair regularization does not beat vanilla SAE on any primary metric at the matched sparsity budgets.}
\label{fig:ruleA}
\end{figure*}

\begin{figure*}[t]
\centering
\includegraphics[width=0.9\linewidth]{main_metrics_ruleB.pdf}
\caption{Seed-level Shapes3D results under rule B, the dense centroid edit. The ranking is similar to rule A, indicating that the negative result is not caused only by a brittle sparse editing heuristic.}
\label{fig:ruleB}
\end{figure*}

\begin{figure}[t]
\centering
\includegraphics[width=\linewidth]{recon_tradeoff.pdf}
\caption{Reconstruction-selectivity trade-off. Increasing the active-latent budget improves selectivity for all methods, but pair regularization does not produce a better frontier point than vanilla SAE.}
\label{fig:tradeoff}
\end{figure}

\subsection{Ablations}

The ablations further weaken the case for the proposed objective. Removing the entire pair-regularization block (\texttt{no\_pair\_regularization}) actually improves the rule-A selective score from $-0.1971$ to $-0.1914$ and raises target success from 0.0260 to 0.0285. Removing concentration alone also slightly improves selective score to $-0.1965$, while removing alignment and separation hurts selective score to $-0.1996$ but raises consistency to 0.5710. These patterns suggest that the proposed terms do change the latent geometry, but not in a way that consistently improves the intervention objective.

\begin{table}[t]
\caption{Budget-B ablations, all at seed 11. The full pair objective does not outperform its own ablations on the primary rule-A metrics.}
\label{tab:ablation}
\centering
\small
\begin{tabular}{lcccc}
\toprule
Method & Recon. $\downarrow$ & Target A $\uparrow$ & Selective A $\uparrow$ & Consistency A $\uparrow$ \\
\midrule
Full pair SAE & 0.3265 & 0.0260 & -0.1965 & 0.5718 \\
No concentration & 0.3261 & 0.0265 & -0.1965 & 0.5653 \\
No align/separation & 0.3263 & 0.0253 & -0.1996 & 0.5710 \\
No pair regularization & \textbf{0.3256} & \textbf{0.0285} & \textbf{-0.1914} & 0.5695 \\
\bottomrule
\end{tabular}
\end{table}

\begin{figure}[t]
\centering
\includegraphics[width=\linewidth]{ablations.pdf}
\caption{Budget-B ablations. Removing the pair losses does not reduce performance on the main selective intervention metric, and in this run it improves it.}
\label{fig:ablations}
\end{figure}

\subsection{Counterfactual retrieval analysis}

Counterfactual retrieval tells the same story as the intervention metrics. The pair model does not improve the consistency rate over vanilla SAE at the higher matched budget, and the dense-rule results closely track the sparse-rule results. Figure~\ref{fig:retrieval} summarizes the cosine-distance comparison against true same-factor counterfactuals and distractors. The one-seed MP-SAE run performs best on this metric family, with rule-A consistency 0.5935 versus 0.5670 for vanilla and 0.5666 for the pair model.

\begin{figure}[t]
\centering
\includegraphics[width=\linewidth]{counterfactual_retrieval.pdf}
\caption{Counterfactual retrieval analysis. The pair-regularized model does not open a clear margin over the matched baselines, while MP-SAE produces the strongest retrieval consistency in the single-seed confirmatory run.}
\label{fig:retrieval}
\end{figure}

\section{Discussion and Limitations}
\label{sec:discussion}

The main limitation is scope. This is a synthetic mechanism study on Shapes3D with one frozen backbone, one selected representation, width 1024, and two sparsity budgets. CelebA transfer was not run in the rerun, so the paper makes no claim about real-world attribute editing. The negative result is therefore local: this specific pair-regularization recipe did not improve intervention outcomes in this specific frozen-feature regime.

The nuisance audit is also important substantively. Even the strict photometric setting falls to 75.6\% agreement on its weakest factor, far below the 99\% threshold. That failure is useful: it shows why nuisance invariance should be validated rather than assumed for frozen self-supervised features. In our rerun, the effective invariance weight becomes zero, so the study should not be read as evidence for or against invariance regularization itself.

Another limitation is that the sparsity-control protocol had to be revised from pure $\ell_1$ control to an explicit active-set cap in order to hit the registered matched-budget windows. We report that change directly because it affects interpretation: the comparison is genuinely matched in realized active-latent count, but it is not a pure-L1-only study. Finally, the MP-SAE comparison uses a single seed, so it supports only a confirmatory statement that stronger sparse coding may matter more than the proposed pair losses.

\section{Conclusion}

We asked whether pair-supervised regularization can improve the selectivity of counterfactual edits in frozen vision SAEs. On Shapes3D with exact single-factor pairs, the answer from this rerun is no: the pair-regularized model slightly improves concentration diagnostics but does not beat both vanilla SAE and Archetypal SAE on the primary intervention metrics at matched sparsity budgets. At budget B under rule A, vanilla SAE achieves the best selective score among the three main methods ($-0.1935 \pm 0.0022$), while the pair-regularized model scores $-0.1971 \pm 0.0007$. A one-seed MP-SAE baseline further improves reconstruction and retrieval consistency, suggesting that the sparse-coding family is a stronger lever than this form of pair regularization in the evaluated regime.

This is a negative result, but it is still useful. It narrows the novelty space for supervised vision SAEs, validates the importance of audit-gated nuisance assumptions, and shows that better latent concentration alone is not enough to claim better causal edit behavior. Future work should test whether stronger paired supervision helps in regimes with verified nuisance pairs, richer data, or sparse-coding families that already dominate the reconstruction-preservation trade-off.

\bibliographystyle{plainnat}
\bibliography{references}

\end{document}
